\documentclass[12pt,twoside]{article}

\input{macros-sp26}
\newcommand{\thelecturenum}{5}

\title{CS336 Lecture 5}

\begin{document}

\handout{Lecture \thelecturenum: GPUs (and TPUs)}

\setlength{\parindent}{0pt}

\subsection*{Engineering the Graphics Processing Unit}

\begin{center}
    \textbf{Goal:} Make GPUs (and TPUs) less magic.
\end{center}

Understand when GPUs get slow. \hfill Understand how to make fast algorithms.

\begin{figure}[hpbt]
    \centering
    \includegraphics[width=1\linewidth]{overview.png}
\end{figure}

\subsubsection*{Before we start..}

Substantial credit goes to a few sources I'd like to highlight ..

\begin{figure}[hpbt]
    \centering
    \includegraphics[width=1\linewidth]{goku.png}
\end{figure}

Horace He's blog (!) \hfill GPU Mode group (!) \hfill TPU (and now GPU) book (!!)

\subsubsection*{Organization today:}

\begin{center}
    \textbf{Part 1}: GPUs in depth - how they work and important parts.
\end{center}

\begin{center}
    \textbf{Part 2}: Understanding GPU performance.
\end{center}

\begin{center}
    \textbf{Part 3}: Putting it together - unpacking FlashAttention.
\end{center}

\subsubsection*{Setting the stage: compute leads to predictable performance}

Often times, compute leads to predictable accuracy gains for language models.

\begin{figure}[hpbt]
    \centering
    \includegraphics[width=0.6\linewidth]{setting.png}
\end{figure}

Faster hardware, better utilization, improved parallelization can drive progress (for now). 

\subsubsection*{How do we get compute scaling? Early on - Dennard scaling}

\begin{figure}[hpbt]
    \centering
    \includegraphics[width=0.6\linewidth]{dennard.png}
\end{figure}

But the traditional form of scaling (\textit{Dennard scaling}) from 1980-2000s has tapped out ... how do we feed LLMs' insatiable appetite for compute?  

\subsubsection*{Parallel scaling continues}

Parallel scaling with GPUs has scaled \(> 1000 \times\) in 10 years.

\begin{claim}
    There is no LLM scaling without GPU scaling.
\end{claim}

\subsubsection*{How is a CPU different from a GPU?}

CPUs optimize for few fast threads WHILE GPUs optimize for many many threads.

\begin{figure}[hpbt]
    \centering
    \includegraphics[width=1\linewidth]{cpu.png}
\end{figure}

Many tiny ALUs. \hfill GPUs optimize for throughput.

Much less support for branching. \hfill CPUs optimize for latency.

\subsubsection*{Anatomy of a GPU (execution units)}

\begin{figure}[hpbt]
    \centering
    \includegraphics[width=0.75\linewidth]{sms.png}
\end{figure}

\begin{itemize}
    \item GPUs have many \textbf{SMs (streaming multiprocessors)} that independently execute "blocks" (jobs).
    \item Each SM further contains many \textbf{SPs (streaming processors)} that can execute "threads" in parallel.
\end{itemize}

\subsubsection*{Anatomy of a GPU (memory)}

\textbf{The closer the memory to the SM, the faster it is} - L1 cache and shared memory is \textit{inside} the SM. L2 cache is on the die and global memory are memory chips next to the GPU. 

\begin{figure}[hpbt]
    \centering
    \includegraphics[width=0.9\linewidth]{bet.png}
\end{figure}

SRAM (shared/cache memory) is much more expensive (100\(\times\)) but about 8\(\times\) faster than DRAM (global memory).

\subsubsection*{Execution model of a GPU}

\begin{figure}[hpbt]
    \centering
    \includegraphics[width=0.9\linewidth]{exe.png}
\end{figure}

There are 3 important players in this execution model.
\begin{itemize}
    \item \textbf{Theads:} Threads "do the work" in parallel - all threads execute the same instructions but with different inputs (SIMT).
    \item \textbf{Blocks:} Blocks are groups of threads. Each block runs on a SM with its own shared memory.
    \item \textbf{Warp:} Threads always execute in a "warp" of 32 consecutively numbered threads each.
\end{itemize}

\subsubsection*{Memory model of a GPU}

Each thread can access its own registers and shared memory within the block.

\begin{figure}[hpbt]
    \centering
    \includegraphics[width=1\linewidth]{mem_mod.png}
\end{figure}

\textbf{Information that goes across blocks needs to be read/written to global memory (slow).}

\subsubsection*{Side thread - What about TPUs?}

GPUs, TPUs, and many other accelerators are at a high level, similar. 

\begin{figure}[hpbt]
    \centering
    \includegraphics[width=1\linewidth]{tpu.png}
\end{figure}

\begin{itemize}
    \item \textbf{Core structure} - lightweight control, fast (big) matmul unit, fast memory.
    \begin{figure}[hpbt]
        \centering
        \includegraphics[width=1\linewidth]{terms.png}
    \end{figure}
    \item Execution difference - no warps (just blocks - tradeoffs in matMul vs nonMatMul).
    \item Networking difference - how accelerators are networked (in parallelism lecture). 
\end{itemize}

\subsubsection*{Strengths of the GPU model}

\begin{itemize}
    \item Easily scales up hard workloads (by adding more SMs).
    \item Easy (?) to program due to the SIMT model.
    \item Threads are "lightweight" and can be stopped/resumed.
\end{itemize}

\subsubsection*{GPUs as fast matrix multipliers}

\begin{itemize}
    \item Early days of NVIDIA GPUs - programmable shaders.
    \item Researchers hacked this to do matmuls.
    \item Volta days of NVIDIA GPUs - new matmul hardware (means matmuls are fast and special).
    \item Tensor cores (introduced in V and T series) are specialized matrix multiplication circuits.
\end{itemize}

\begin{center}
    \textbf{Matmuls are \(> 10\times\) faster than other floating point ops!}
\end{center}

\subsubsection*{GPUs - what are they and how do they work}

\begin{center}
    GPUs are massively parallel - same instructions applied across many workers.
\end{center}

\begin{center}
    Compute (and especially matmuls) have scaled faster than memory.
\end{center}

\begin{center}
    We have to respect the memory hierarchy to make things go fast.
\end{center}

\subsection*{Making ML workloads fast on a GPU}

Performance on a GPU can be complex, even for something as simple as a square matmul. 

\begin{figure}[hpbt]
    \centering
    \includegraphics[width=1\linewidth]{roofline_gpu_vs_cpu.png}
\end{figure}

Key to this section: \textbf{how do we avoid being memory bound?}

\begin{enumerate}
    \item Control divergence (not a memory bottleneck)
    \item Low precision computation
    \item Operator fusion
    \item Recomputation
    \item Coalescing memory
    \item Tiling
\end{enumerate}

\subsubsection*{Control divergence (not a memory issue)}

\begin{itemize}
    \item GPUs operate in a SIMT model - every thread in a warp is executing the same instruction.
    \item Conditionals are fine, but lead to significant overhead from the execution model. 
\end{itemize}

\subsubsection*{Trick 1: Low precision computation}

If you have fewer bits, you have fewer bits to move.

\begin{center}
    \textbf{Example:} elementwise ReLU (\(x = \max(0, x)\) on a vector of size \(n\).
\end{center}

(Float 32 case)
\begin{itemize}
    \item \textbf{Memory access}: 1 read (\(x\)), 1 write (if \(x < 0\)), fp32 \(=\) 4 bytes per op.
    \item \textbf{Operations}: 1 comparison op, 1 FLOP.
    \item \textbf{Intensity Inverse}: 8 bytes / FLOPs.
\end{itemize}

(Float 16 case)
\begin{itemize}
    \item \textbf{Memory access}: 1 read (\(x\)), 1 write (if \(x < 0\)), fp16 \(=\) 2 bytes per op.
    \item \textbf{Operations}: 1 comparison op, 1 FLOP.
    \item \textbf{Intensity Inverse}: 4 bytes / FLOPs.
\end{itemize}

\subsubsection*{Low precision drives faster matrix multiples}

Lots of operations in modern GPUs are accelerated via low / mixed precision operations.

\begin{figure}[hpbt]
    \centering
    \includegraphics[width=0.75\linewidth]{tensor_cores.png}
\end{figure}

Operations that can use 16-bit storage (fp16/bf16)
\begin{itemize}
    \item Matrix multiplication.
    \item Most pointwise operations (e.g. relu, tanh, add, sub, mul).
\end{itemize}

Operations that need ore precision (fp32/fp16)
\begin{itemize}
    \item Adding small values to large sums can lead to rounding errors.
    \item Reduction operations (e.g. sum, softmax, normalization).
\end{itemize}

Operations that need more range (fp32/bf16)
\begin{itemize}
    \item Pointwise operations where \(\vert f(x) \vert >> \vert x \vert\) (e.g. exp, log, pow).
    \item Loss functions.
\end{itemize}

\subsubsection*{Frontier in low precision}

Very low precision (FP8) with different tradeoffs. \hfill Multiple scaling factors MXFP8 (Blackwell).

\begin{figure}[hpbt]
    \centering
    \includegraphics[width=1\linewidth]{fp8_mxfp8.png}
\end{figure}

\subsubsection*{Trick 2: Operator fusion}

Think of a GPU like a factory - inputs come from a warehouse (memory) and is processed at a factory.

\begin{itemize}
    \item What if we have to do many operations? Shipping back and forth is somewhat silly.
    \item \textbf{Example} - Computing \(\sin^2 x\ + \cos^2 x\) naively launches 5 CUDA kernels (back and forth).
\end{itemize}

All 5 pointwise operations can be fused into a single CUDA kernel call. "Easy" fusions like this can be donw automatically by compilers (torch.compile).

\subsubsection*{Trick 3: Recomputation}

\begin{figure}[H]
    \centering
    \includegraphics[width=1\linewidth]{backprop.png}
\end{figure}

In backpropagtion, we store the activations (yellow) and compute Jacobians (green). 

\subsubsection*{Storing (and retrieving) activations can be expensive!}

Let's say we stack 3 sigmoids on top of each other.

\begin{figure}[hpbt]
    \centering
    \includegraphics[width=0.5\linewidth]{bad.png}
\end{figure}

This is really terrible for performance - 8 memory read/writes, very low arithmetic intensity.

\subsubsection*{Throw away the activations, re-compute them!}

\begin{figure}[hpbt]
    \centering
    \includegraphics[width=0.5\linewidth]{recom.png}
\end{figure}

Throwing away computation can be actually be optimal, with \(5 / 8\)th the memory accesses!

\subsubsection*{Trick (?) 4: DRAM and Memory coalescing}

\textbf{DRAM} (global memory) is read in "burst mode" - each read gives you many bytes!

\begin{figure}[hpbt]
    \centering
    \includegraphics[width=1\linewidth]{burst_sections.png}
\end{figure}

Burst mode comes from the slow per-row copy to the sense amplifier.

\subsubsection*{Memory coalescing}

Memory accesses are \textit{coalesced} if all the threads (in a warp) fall within the same burst.

\begin{figure}[hpbt]
    \centering
    \includegraphics[width=1\linewidth]{coalesced_loads.png}
\end{figure}

\subsubsection*{Coalescing for matrix multiplication}

\begin{figure}[hpbt]
    \centering
    \includegraphics[width=1\linewidth]{coalesced_access.png}
\end{figure}

For row-major matrices - \textbf{threads that move along rows are not coalesced}.

\begin{remark}
    Note how the second diagram reads the entire vector at each step!
\end{remark}

\subsubsection*{Trick 5 (the big one): Tiling}

\textbf{Tiling} is the idea of grouping and ordering threads to minimize global memory access.

\begin{center}
    Let's go back to matrix multiplication .. Note that memory access is not coalesced and repeated
\end{center}

\subsubsection*{Tiling - store and reuse information in shared memory}

Cut up the matrix into smaller "tiles" and load this into shared memory.

\begin{figure}[hpbt]
    \centering
    \includegraphics[width=1\linewidth]{tiled_matmul.png}
\end{figure}

\begin{center}
    \textbf{Advantages}: repeated reads now access shared, not global memory and memory access can be coalesced.
\end{center}

\subsubsection*{Tiling math}

\textbf{Non-tiled matrix multiply}: each input is read \(N\) times from global memory.

\begin{figure}[hpbt]
    \centering
    \includegraphics[width=1\linewidth]{tiled_loops.png}
\end{figure}

\textbf{Tiled matrix multiply}: each input is read \(N / T\) times from global memory and \(T\) times within each tile. This is a factor of \(T\) reduction in global memory access

\subsubsection*{Complexities with tiling}

\textbf{Tile sizes may not divide the matrix dimensions} and lead to low utilization.

\begin{figure}[hpbt]
    \centering
    \includegraphics[width=1\linewidth]{tile_quantization.png}
\end{figure}

Factors affecting tile sizes:
\begin{itemize}
    \item Coalesced memory access.
    \item Shared memory size.
    \item Divisibility of the matrix dimension. 
\end{itemize}

\subsubsection*{Complexities with tiling 2 - memory alignment}

Loading tiles are fast if bursts align with the matrix.

\begin{figure}[hpbt]
    \centering
    \includegraphics[width=0.9\linewidth]{aligment.png}
\end{figure}

\textbf{Coalesced accesses may be impossible depending on the dimension of the matrix ..} (have to do padding).

\subsubsection*{Putting it together: understanding a matrix mystery}

\begin{center}
    Why is it \textit{faster} to have bigger matrices?
\end{center}

We understand some of this (compute intensity, tiling blocks). Let's take a closer look ..

\begin{itemize}
    \item Tiling has a major impact through alignment.
    \item Quantization has a nontrivial impact on periodic behavior.
\end{itemize}

\begin{center}
    \textbf{An A100 has 108 SMs, so it cannot execute all 120.}
\end{center}

\subsection*{Understanding Flash Attention}

\textbf{Flash attention} [Dao et al] dramatically accelerates attention .. but how?

\begin{figure}[hpbt]
    \centering
    \includegraphics[width=1\linewidth]{flashattention.png}
\end{figure}

\textbf{Technique from paper}: We apply two established techniques (tiling, recomputation) to overcome the technical challenge of computing exact attention in sub-quadratic HBM accesses. 

\subsubsection*{Part 1: tiling for the KQV matrix multiply}

\begin{figure}[hpbt]
    \centering
    \includegraphics[width=0.6\linewidth]{flashattention_diagram.png}
\end{figure}

\begin{center}
    This figure from the paper is literally just tiling for a KQV matrix multiply .. \textbf{but what do we do about the softmax?}
\end{center}

\subsubsection*{Part 2: incremental computation of the softmax}

From Mikailov and Gimelshein 2018,

\begin{figure}[hpbt]
    \centering
    \includegraphics[width=1\linewidth]{safe_softmax_algorithm.png}
\end{figure}

To keep track of the max, incrementally update the max and set up a telescoping sum.

\begin{center}
    \textbf{This lets you compute the softmax tile-by-tile}.
\end{center}

\subsubsection*{Putting it all together - the forward pass of flash attention}

\begin{figure}[hpbt]
    \centering
    \includegraphics[width=1\linewidth]{flashattention_blocks.png}
\end{figure}

From Dao 2023, we see
\begin{itemize}
    \item Tile-wise computation of the inner products (\(S\)).
    \item Fusion of the exponential operator.
    \item Tile-wise computation of the softmax via the online, telescoping sum trick.
\end{itemize}

\begin{center}
    (We won't cover the backward pass - but they recompute tile-by-tile ..)
\end{center}

\subsubsection*{Recapping the Graphics Processing Unit}

\begin{center}
    Hardware powers scale and low-level details determine what scales or doesn't.
\end{center}

\begin{center}
    Current GPU-based compute strongly encourages thinking about matmul \(+\) data movement.
\end{center}

\begin{center}
    Thinking about the GPU (coalescing, tiling, fusion) leads us to good performance. 
\end{center}

\end{document}
